%% kapa: exam=1 round=23 year=2012 period=1 form=A subject=economics subject_ko=경제원론 seq=31 number=71
%% source: Q-Net 감정평가사 기출문제 (cst00309) artlSeq=5052578 file=raw/1008/5052578/2012년 제23회 감정평가사 제1차시험 기출문제(1교시 A형).pdf
%% conversion: kordoc 4.24.51 -> tools/gen1.py (auto); review=proofread-sonnet; flags=none
\begin{question}{71}
A국의 연간 실질 GDP 변화율과 실업률의 변화가 다음과 같은 관계에 있다. 이에 관한 설명으로 옳은 것은?
\begin{framedtext}
실질 GDP 변화율(\%) = 3\% - 2 × 실업률(\%)의 변화
\end{framedtext}
\begin{choices}
  \item 자연실업률은 3\%이다.
  \item 실업률이 5\%에서 6\%로 상승하면 실질 GDP는 2\% 감소한다.
  \item 실질 GDP가 1\% 하락하면 실업률은 5\%에서 5.5\%로 상승한다.
  \item 물가가 상승하면 단기적으로 실질 GDP가 감소한다.
  \item 실업률이 변화하지 않을 경우 실질 GDP는 3\% 증가한다.
\end{choices}
\end{question}
\answer{5}
